\documentclass[12pt, letterpaper]{article}
\setlength{\parindent}{0pt}

\usepackage{xcolor}
\usepackage[bottom=0.5cm, top=0.5cm, left=1.5cm, right=1.5cm]{geometry}
\usepackage{hyperref}
\usepackage{amsmath}

\hypersetup{
    colorlinks = true,
    linkcolor = blue,
    filecolor = magenta,
    urlcolor = cyan,
}

\title{MAS3114}
\author{Ian Zou}

\begin{document}

\maketitle

\part*{Unit 1 - Linear Systems \& Matrices}

\section*{Lecture 1}

\subsection*{Linear Equations}

A \textbf{linear equation} is in the form:

\medskip
{\centering

$ a_1 x_1 + a_2 x_2 + . . . + a_n x_n = b_1 $

}
\medskip

Where $ x_1, x_2, ..., x_n $ are variables, $ a_1, a_2, ..., a_n $ and $ b $ are scalars, and $ n $ is a positive integer.

\medskip

Examples of linear equations include:
\begin{itemize}
    \item $ 5x_1 + 4x_2 = 0 $
    \item $ \sqrt{5}x_1 + 4x_2 = 0 $
\end{itemize}

The following are \textit{NOT} linear equations:
\begin{itemize}
    \item $ 5\sqrt{x_1} + 4x_2 = 0 $
    \item $ 5\sqrt{x_1} + 4x_2 = 0 $
\end{itemize}

\subsection*{Linear Systems}

A \textbf{system of linear equations} (\textbf{linear system}) is a collection of one or more linear equations.
\medskip

The following is a linear system:

\medskip
{\centering

$ 2x_1 + x_2 + x_3 = 1 $
$ x_1 - 2x_2 - x_3 = 2 $

}
\medskip

A linear system is:

\begin{itemize}
    \item \textbf{Consistent}: has 1 solution or \textit{infinitely many} solutions.
    \item \textbf{Inconsistent}: has \textit{no} solutions.
\end{itemize}

Two linear systems are \textbf{equivalent} if they have the same solution set.
\medskip

\subsection*{Matrix Notation}

Used to represent a linear system.
\medskip

The following is an example of a $ 3 \times 3 $ \textbf{coefficient matrix}:

\medskip
{\centering

$ A = \begin{bmatrix}
    a_{11} & a_{12} & a_{13} \\
    a_{21} & a_{22} & a_{23} \\
    a_{31} & a_{32} & a_{33}
\end{bmatrix}  $

}
\medskip

The \textbf{size} of a matrix is the number of rows and columns. Any $ m \times n $ matrix is a 2-D array of numbers.

\begin{enumerate}
    \item \textit{Example}: Solve the system.

        {\centering

            $ x_1 - x_2 = -4 $

            $ x_1 + 2x_2 = 2 $

        }

        \bigskip

        \textit{Naive Solution}: We get the singular equation by subtracting the equations: % TODO: clarify 

        {\centering

            $ 3x_2 = 6 $

            $ x_2 = 2 $

        }

        \medskip

        We can plug it back into one of the equations to solve for $ x_1 $:

        {\centering

            $ x_1 + 2(2) = 2 $

            $ x_1 = -2 $

        }

        \medskip

        Thus, the solution set is:

        {\centering

            $\begin{bmatrix}
                -2 \\
                 2 \\
            \end{bmatrix}$

        }

\end{enumerate}

\subsection*{Elementary Row Operations}

\begin{enumerate}
    \item \textbf{Interchange}: Interchange two rows.
    \item \textbf{Scaling}: Multiply every entry in a row by a nonzero constant.
    \item \textbf{Replacement/Row Addition}: Add a multiple of one row to another row.
\end{enumerate}

Two matrices are \textbf{row equivalent} if there is a sequence of elementary row operations that transform one matrix into the other.
\medskip

\subsection*{Echelon Forms}

% TODO: echelon forms

\end{document}
